% Exact finite obstruction for the literal descent investigation.
% This fragment uses only the actual source quotient described below.
\subsection{A canonical central refinement of the finite descent source}

Put $B_0=\langle s,r\mid s^2=r^3=1\rangle$, $\tau=sr$,
$x=\tau^{-2}$ and $y=sxs$.  Let $F_0=\langle x,y\rangle$ be
the index-six free subgroup, with its specified positive peripheral
marking.  Write
\[
 S=\begin{pmatrix}0&1\\-1&0\end{pmatrix},\qquad
 U=\begin{pmatrix}1&1\\0&1\end{pmatrix},\qquad
 R=\begin{pmatrix}0&-1\\1&-1\end{pmatrix}.
\]
The two maps used to construct the finite source are
\[
\begin{aligned}
 \theta&:F_0\longrightarrow\operatorname{SL}_2(\mathbb Z/12),
 \quad (x,y)\longmapsto(S,U^{-1}),\\
 q&:B_0\longrightarrow
 Q:=\operatorname{SL}_2(\mathbb Z/12)/\{\pm I\},
 \quad(s,r)\longmapsto([S],[R]).
\end{aligned}                                  \tag{B.M1}
\]
Define the actual normal subgroup and quotient
\[
 K=\ker q\cap\operatorname{core}_{B_0}(\ker\theta),
 \qquad D=B_0/K.                               \tag{B.M2}
\]
The core is the intersection of the six conjugates indexed by
$1,s,r,r^2,sr,sr^2$.  Thus (B.M2) is a finite quotient of the
specified discrete source, and also a continuous quotient of
$\widehat B_0$.  It is exactly the degree-$168$ permutation
representation in \texttt{check\_descent\_source\_frames.g}:
the six induced degree-$20$ blocks are the actions on primitive
vectors modulo four and nonzero vectors modulo three; the final
degree-$48$ block is the action of $Q$ on primitive vectors modulo
twelve modulo their common sign.

\begin{result}{B.2.6}{Proved}{The canonical relative central obstruction}
Let $\overline s,\overline r$ be the images of $s,r$ in the actual
group $D$ of (B.M2).  There is a finite, canonically marked central
extension
\[
 1\longrightarrow M\longrightarrow D^{\mathrm c}
   \xrightarrow{\pi}D\longrightarrow1,
 \qquad
 D^{\mathrm c}=B_0/[B_0,K],\quad M=K/[B_0,K].    \tag{B.M3}
\]
Write $S_0,R_0$ for its distinguished images of $s,r$.
For every actual automorphism of $D$ of the form
\[
 \alpha_f(\overline s)=\overline s f_D,
 \qquad
 \alpha_f(\overline r)=f_D^{-1}\overline r,
 \qquad f_D\in D,                              \tag{B.M4}
\]
the following element of the finite group $M$ is well-defined,
independently of the choice of a lift $\widetilde f$ of $f_D$:
\[
 a=(S_0\widetilde f)^2,\qquad
 b=(\widetilde f^{-1}R_0)^3,\qquad
 o(\alpha_f)=a^3b^2.                           \tag{B.M5}
\]
Let $\ell$ be the map from automorphisms of $D^{\mathrm c}$
preserving $M$ and fixing $S_0R_0$ to the automorphisms of
$D$ in (B.M4).  Its exact image and fibres are
\[
              \operatorname{im}\ell=o^{-1}(1),
              \qquad |\ell^{-1}(\alpha_f)|=1
                  \quad(\alpha_f\in o^{-1}(1)).
\]
On this image the unique automorphism is given by
\[
 f^{\mathrm c}=\widetilde f\,ab,
 \qquad S_0\longmapsto S_0f^{\mathrm c},
 \qquad R_0\longmapsto(f^{\mathrm c})^{-1}R_0.   \tag{B.M6}
\]
This is a finite, effectively computable necessary obstruction
for a normalized gentle automorphism inducing (B.M4).
In addition, every such gentle lift has
$f^{\mathrm c}\in P'$, where $P$ is the image of $F_0$ in
$D^{\mathrm c}$, and preserves the positive conjugacy classes
of $S_0$ and $R_0$.  Prescribed modular residues of the two
conjugators give the exact finite tests in (B.M8) below.
\end{result}

\begin{proof}
First the map $B_0^{\mathrm{ab}}\to D^{\mathrm{ab}}$ is an
isomorphism, with both groups $C_2\times C_3$.  Indeed $q$
factors through $D$.  Reduction of $Q$ modulo two, followed by
the sign on $\operatorname{SL}_2(\mathbb F_2)\cong S_3$, detects
the $C_2$ generated by $s$.  Reduction modulo three, followed
by $\operatorname{SL}_2(\mathbb F_3)\to C_3$, detects the $C_3$
generated by $r$.  Here the kernel of the latter map is the
normal quaternion group consisting of $\pm I$ and the six
trace-zero matrices; $S$ belongs to this kernel and the
order-three matrix $R$ does not.  Both maps kill the common
sign, so they are maps from the displayed $Q$.  They give a
surjection $D\to C_2\times C_3$ inverse on abelianization to
the quotient of $B_0$.  In particular $K\subseteq B_0'$.

The group $M$ in (B.M3) is central and finitely generated
abelian: $K$ has finite index in the finitely generated group
$B_0$, and its image is central in $D^{\mathrm c}$.  Moreover
$M\subseteq(D^{\mathrm c})'$ since $K\subseteq B_0'$.  The
transfer homomorphism from $D^{\mathrm c}$ to its finite-index
central subgroup $M$ sends $m\in M$ to $m^{|D|}$.  It kills
$(D^{\mathrm c})'$, since its target is abelian.  Consequently
every $m\in M$ satisfies $m^{|D|}=1$.  A finitely generated
abelian group of bounded exponent is finite.  This proves
finiteness of both $M$ and $D^{\mathrm c}$, without replacing
the actual source by an arbitrary covering group.

For completeness, this particular kernel is also the Schur
multiplier of $D$.  Let $L$ be the free group on $s,r$,
$V=\langle\!\langle s^2,r^3\rangle\!\rangle$, and $W$ the
kernel of $L\to D$.  The image of $W$ in $L^{\mathrm{ab}}$
is exactly $2\mathbb Z\oplus3\mathbb Z$, the image of $V$.
The group $V/[L,V]$ is generated by the images of $s^2,r^3$;
their independent images in $L^{\mathrm{ab}}$ show that
$V\cap L'=[L,V]$.  It follows, by directly eliminating these
two abelian generators, that
\[
 M= W/(V[L,W])
   \cong (W\cap L')/[L,W]=H_2(D,\mathbb Z).      \tag{B.M7}
\]
The final equality is the Hopf presentation formula.  The
preceding transfer argument already proves the finiteness
needed here, independently of this identification.

Because (B.M4) respects the two defining relations of $D$,
the elements $a,b$ in (B.M5) belong to $M$.  Replacing
$\widetilde f$ by $\widetilde f z$, with $z\in M$, changes
them to $az^2,bz^{-3}$.  Hence $a^3b^2$ is independent of
the lift.  To impose the two cycle relations upstairs one
must solve $z^2=a^{-1}$ and $z^3=b$.  Those equations force
$z=z^3(z^2)^{-1}=ab$.  Substitution shows that this unique
choice solves both equations exactly when $a^3b^2=1$.
This proves the formula and uniqueness in (B.M6).

The two corrected generator images define a homomorphism
$B_0\to D^{\mathrm c}$, and their product is $S_0R_0$.
Its composite with $\pi$ is $\alpha_f$ composed with the
original quotient.  It therefore takes $K$ into $M$ and
kills $[B_0,K]$.  Thus it gives an endomorphism $\widetilde
\alpha_f$ of $D^{\mathrm c}$.  Put $J=\operatorname{im}
\widetilde\alpha_f$.  Since $\alpha_f$ is onto,
$D^{\mathrm c}=JM$.  Centrality of $M$ then gives
$(D^{\mathrm c})'=J'\subseteq J$.  Since $M$ is contained
in this derived subgroup, $M\subseteq J$, whence
$J=D^{\mathrm c}$.  Finiteness makes the endomorphism an
automorphism.

The open kernel of $\widehat B_0\to D^{\mathrm c}$ is the
closure of $[B_0,K]$, and the inverse image of $1\in D$ is
the closure of $K$.  A continuous automorphism of
$\widehat B_0$ inducing (B.M4) preserves the latter closed
subgroup and its closed commutator with $\widehat B_0$.
It therefore induces the just proved unique automorphism of
$D^{\mathrm c}$ whenever it fixes $\tau$.  For a normalized
gentle pair the element $f$ belongs to the closed derived
subgroup of $\widehat F_0$, and its two torsion images are
conjugate to the original positive generators.  Their
images in this actual finite quotient have the asserted
derived membership and positive conjugacy properties.
\end{proof}

\paragraph{Exact residue fibres for the two positive frames.}
The map $q$ in (B.M1) induces $q_{\mathrm c}:D^{\mathrm c}
\to Q$.  Its kernel $N$ is contained in $P$: reduction
modulo two is the specified quotient $B_0/F_0\cong S_3$.
Take prescribed residues $\beta_Q,A_Q\in q_{\mathrm c}(P)$
and any lifts $\beta_0,A_0\in P$.  There is a positive
$S_0$-conjugator in the fibre $q_{\mathrm c}^{-1}(\beta_Q)
\cap P$, and a positive $R_0$-conjugator in the fibre
$q_{\mathrm c}^{-1}(A_Q)\cap P$, precisely when
\[
\begin{split}
 &\beta_0^{-1}(S_0 f^{\mathrm c})\beta_0
       \text{ is conjugate to }S_0
       \text{ in }\langle N,S_0\rangle,\\
 &A_0^{-1}((f^{\mathrm c})^{-1}R_0)A_0
       \text{ is conjugate to }R_0
       \text{ in }\langle N,R_0\rangle.         \tag{B.M8}
\end{split}
\]
Indeed every possible conjugator is respectively
$\beta_0n$ or $A_0n$ with $n\in N$.  Conjugacy of $S_0$
under $\langle N,S_0\rangle=N\langle S_0\rangle$ is
exactly conjugacy under $N$, because $S_0$ commutes with
its own powers.  The same argument applies to $R_0$.
Thus (B.M8) retains the actual prescribed residue fibres;
unrestricted conjugacy alone would omit these conditions.

\paragraph{The elementary central two-quotient is canonical.}
The finite group $D^{\mathrm c}/M^2$ is the canonical
relative central quotient of this marked source with the
maximal elementary abelian two-kernel.  Its kernel is
$M/2M$.  Central
extensions of $D$ by $\mathbb F_2$ are represented by
$H^2(D,\mathbb F_2)$ with trivial action.  Such an extension
receives the marked map from $B_0$ exactly when its
restriction to $\langle\overline s\rangle$ is zero:
one must lift $s$ to an involution, while the lift of $r$
can always be corrected to order three, since multiplication
by three is invertible on the central two-kernel.
Thus the admissible classes are exactly
\[
 \ker\bigl(H^2(D,\mathbb F_2)\longrightarrow
 H^2(\langle\overline s\rangle,\mathbb F_2)\bigr)
       \cong\operatorname{Hom}(M,\mathbb F_2).  \tag{B.M9}
\]
The isomorphism follows also from (B.M7) and the elementary
abelianization calculation: the single extension-class
direction arising from $D^{\mathrm{ab}}\cong C_6$ is
detected on the source involution.  Taking all admissible
cocycles together therefore constructs the actual marked
quotient $D^{\mathrm c}/M^2$, rather than a selected abstract
cover with unspecified source map.  Formulae (B.M5),
(B.M6) and (B.M8) apply to this quotient verbatim.

\paragraph{Executed calculation for the twenty-four source candidates.}
The scripts \texttt{check\_descent\_central2.g} and
\texttt{check\_descent\_central2\_joint.g}, with their saved
\texttt{.out} files, compute the cocycles using GAP's exact
polycyclic-group cohomology algorithm.  The result is
\[
 \dim_{\mathbb F_2}H^2(D,\mathbb F_2)=4,
 \qquad \dim_{\mathbb F_2}(M/2M)=3,
 \qquad |D^{\mathrm c}/M^2|=3131031158784.        \tag{B.M10}
\]
The script forms one simultaneous extension from a basis
of the three admissible cocycles, verifies that the
distinguished lifts $S_0,R_0$ generate it, and verifies its
central kernel has order eight.  The image of $F_0$ has
order $521838526464$, and its derived subgroup has order
$3623878656$.  Each of the twenty-four source candidates
from the preceding descent calculation has trivial
obstruction (B.M5).  Its unique corrected lift belongs
to that derived subgroup, has both positive torsion
conjugacy classes, and passes both residue-fibre tests
(B.M8).  Twelve candidates have each displayed $\Delta$
bit.  These are conclusions about the specified finite
quotient and actual elements in it.  This central
refinement excludes none of those candidates; it supplies
no full profinite lifting assertion.

\subsection{A finite certificate for a full positive lift}

The full cover in (B.M3), as opposed to one of its proper
central quotients, supplies a complete finite lifting test.
Here the finite transport theorem is used before taking
any nilpotent residual: its proof needs only normality of
the kernel.  This is stated explicitly immediately before
R1, (D.50.17), and independently checked in the supplied
referee report, the rows for (D.50.15)--(D.50.17).

Let $\mathcal A_D$ be the finite set of triples
$(\alpha,\beta,A)$, where $\alpha\in\operatorname{Aut}(D)$
fixes $\overline\tau=\overline s\overline r$ and
\[
 \alpha(\overline s)=\beta\overline s\beta^{-1},\qquad
 \alpha(\overline r)=A\overline rA^{-1},\qquad
 \beta,A\in\overline F_0.                       \tag{B.M11}
\]
Define $\mathcal A_{\mathrm c}$ by the identical conditions
on $D^{\mathrm c}$, retaining its distinguished source
generators and requiring preservation of $M$.  Let
$\mathcal A_{\infty}$ be the set of actual continuous
automorphisms of $\widehat B_0$ fixing $\tau$, preserving
$\overline K$, and having specified pure positive
conjugators $\beta,A\in\widehat F_0$ for $s,r$.
There are literal reduction maps from both sets to
$\mathcal A_D$.

\begin{result}{B.2.7}{Proved}{An exact finite-certificate lifting equality}
For the marked source (B.M1)--(B.M3), the two reduction
maps have the same image:
\[
 \operatorname{im}(\mathcal A_{\infty}\longrightarrow
                         \mathcal A_D)
   =\operatorname{im}(\mathcal A_{\mathrm c}\longrightarrow
                         \mathcal A_D).         \tag{B.M12}
\]
In particular the right side is an effective, finite
description of all the indicated full positive lifts,
including the two prescribed $D$-conjugator residues.
The full gentle parameter is obtained from each such
lift by the exact inner correction (B.M13) below.  Its
induced outer automorphism of every peripheral quotient,
including the literal pants quotient defining $D_P$,
is unchanged by that correction.
\end{result}

\begin{proof}
Reduction of a full lift to $D^{\mathrm c}$ was proved
in B.2.6, including its actual two conjugator coordinates.
This proves the inclusion from left to right.

For the other inclusion, fix an element
$(\alpha^{\mathrm c},\beta^{\mathrm c},A^{\mathrm c})$
of $\mathcal A_{\mathrm c}$, and write its image as
$(\alpha,\beta,A)$.  Work in the compact space
$\widehat B_0^3$ of the proposed conjugators of
$s,r,\tau^{-1}$.  Prescribe their images in $D$ to be
$\beta,A,1$, respectively.  Impose the positive
peripheral product relation and generation in every
finite quotient.  Every finite collection of these
conditions is contained in an actual joint finite
quotient $q_E:\widehat B_0\twoheadrightarrow E$ refining
$D$.

Put $N=\ker(E\to D)$.  Apply exactly the finite
scalar-cover and subgroup-count construction R1,
(D.50.5)--(D.50.16), to the extension
$1\to N\to E\to D\to1$.  The counting argument requires
only that $N$ be normal.  The separate equality
$[E,N]=N$ is used only in (D.50.17), which is not needed
here.  The construction provides an actual marked
source quotient
\[
 q^\dagger:\widehat B_0\twoheadrightarrow D^\dagger,
 \qquad D^\dagger\twoheadrightarrow D
\]
whose kernel is central, with all the subgroup counts
and actual auxiliary source maps retained.  Because
$q^\dagger(\overline K)$ is central, this map kills
$[\widehat B_0,\overline K]$.  It consequently factors
through the same finite marked group $D^{\mathrm c}$:
$q^\dagger=vq_{\mathrm c}$.

Define the new auxiliary source map by
$q_1^\dagger=v\alpha^{\mathrm c}q_{\mathrm c}$.  This is
an epimorphism.  It has the three positive conjugators
$v(\beta^{\mathrm c}),v(A^{\mathrm c}),1$ relative to
the original source map $q^\dagger$, and its map to $D$
is $\alpha q_D$.  No automorphism of $D^\dagger$ is
asserted or required: $\alpha^{\mathrm c}$ need not
preserve $\ker v$.

The exact equality of epimorphism counts in (D.50.16)
now supplies an epimorphism $\widehat B_0\to E$ whose
positive conjugators have the prescribed three
$D$-values.  The reference count is positive because
$q_E$ itself occurs in it.  Thus every finite
collection of the conditions in the compact space is
satisfied.  Compactness produces full conjugators
$b_s,b_r,b_\tau$ and a continuous endomorphism $\Psi$
of $\widehat B_0$ inducing $\alpha$ on $D$ and
generating every finite quotient.  Its image is compact
and hence closed, so it is onto.  A finitely generated
profinite group is Hopfian: inverse image permutes the
finite set of open subgroups of each index under an
epimorphism, forcing its kernel into every open
subgroup.  Thus $\Psi$ is an automorphism.

The prescribed $D$-value of $b_\tau$ is one, so
$b_\tau\in\overline K$.  Conjugating $\Psi$ by
$b_\tau^{-1}$ fixes $\tau$ literally and leaves its
$D$-action and its two prescribed $D$-conjugator
residues unchanged.  The resulting conjugators are
pure because their $D$-values are pure and $D\to S_3$
is the specified source permutation quotient.  The
resulting automorphism therefore belongs to
$\mathcal A_{\infty}$, proving (B.M12).
\end{proof}

\paragraph{The exact commutator normalization.}
Let $\Phi\in\mathcal A_{\infty}$ and put $f=s\Phi(s)$.
Its action is $s\mapsto sf$, $r\mapsto f^{-1}r$ and
$\tau\mapsto\tau$.  Its restriction to the pure free
group fixes $x$ and sends $y$ to $f^{-1}yf$.
The two actual torsion relations give the two gentle
cycle identities.  In particular, writing the pure
abelianization of $f$ as $(a,b)$, the first cycle gives
$b=-a$.  The exact correction
\[
 \Phi_0=\operatorname{Inn}(x^a)\Phi,\qquad
 f_0=y^a f x^{-a},\qquad f_0^{\mathrm{ab}}=0      \tag{B.M13}
\]
fixes $\tau$ and gives the normalized full multiplier-one
gentle pair $(1,f_0)$.  It changes the literal $D$-action
by $\operatorname{Inn}(\overline x^{a})$; this change
is retained explicitly.  On the actual pants quotient
$\varphi_P:\widehat F_0\to\widehat B_0$, it changes
the descended automorphism by
$\operatorname{Inn}(\varphi_P(x)^a)
=\operatorname{Inn}(s^a)$.  Thus its descended outer
class, and every character of that class, is exactly
the class supplied by the finite-certificate lift.

The normalized multiplier-one gentle representation in
$\operatorname{Out}(\widehat B_0)$ is faithful with
these markings.  Indeed suppose two normalized
representatives differ by $\operatorname{Inn}(b)$.
Their restrictions to the characteristic pure group
both act identically on $\widehat F_0^{\mathrm{ab}}$.
Hence the deck action of $b\bmod\widehat F_0$ on this
abelianization is trivial.  The displayed cyclic and
swap matrices give a faithful copy of $S_3$ on
$\widehat{\mathbb Z}^{2}$, so $b\in\widehat F_0$.
Both representatives fix $x$, and the primitive
centralizer calculation gives $b=x^t$.  If their
commutator pair coordinates are $f_1,f_2$, equality
of the $y$-images gives
$f_2=y^c f_1x^{-t}$ for some $c\in\widehat{\mathbb Z}$.
Taking abelianizations gives $(-t,c)=(0,0)$, hence
$b=1$ and $f_1=f_2$.  This proves the claimed
faithfulness by the actual inner comparison, and
justifies passing from an unchanged descended outer
class to its unique normalized gentle parameter.
